\section{Заключение}

\subsection{Итоги исследования}\label{ux438ux442ux43eux433ux438-ux438ux441ux441ux43bux435ux434ux43eux432ux430ux43dux438ux44f}

Курсовая работа отвечает на продуктовый вопрос: когда в онлайн-эксперименте по CTR целесообразно адаптивное назначение трафика (MAB), а когда --- фиксированное распределение A/B, и как корректно интерпретировать логи для проверки гипотезы о превосходстве одного варианта над другим.

\textbf{Польза адаптивного эксперимента.} На контролируемой синтетике и в попарной постановке с двумя вариантами (E1, E6, E8b, E12) Thompson Sampling и $\varepsilon$-greedy накапливают больше кликов и меньше regret, чем равномерное разбиение \texttt{fixed\_ab}. При пяти руках и горизонте 5000 шагов недобор кликов относительно оракула у TS примерно на 14--16 пунктов меньше, чем у \texttt{fixed\_ab}; при 20 независимых прогонах это преимущество не сводится к одному удачному seed. В постановке «вариант B против A» TS даёт около 140 дополнительных кликов на 10 000 показов. Контекстный LinUCB при информативных признаках в среднем выигрывает у \texttt{fixed\_ab} по regret (26,6 против 49,2 на синтетике E8b), но разброс между прогонами велик --- строго доказать устойчивое превосходство не удалось. \emph{Итог:} бандиты оправданы, когда важно не терять клики во время теста; накопленный CTR бандита сам по себе не доказывает, что B статистически лучше A.

\textbf{Статистический вывод по логам.} Если после bandit-эксперимента применить обычный z-критерий так, будто трафик всё время делили поровну, вывод искажается. Когда варианты A и B на самом деле одинаковы, тест ошибочно заявляет об эффекте в 14,5\% прогонов (E4) вместо допустимых 5\%; когда эффект есть, но слабый (+1 п.п. к CTR), он обнаруживается лишь в 55\% случаев против 85\% у корректного A/B с \texttt{fixed\_ab}. Причина --- неравномерное распределение показов между группами после адаптивного назначения (иллюстрация E5). IPS-взвешенный статистический вывод на bandit-логах возвращает контроль ошибки I рода к 5\%, однако при слабом эффекте мощность остаётся низкой (24\% при +1 п.п., 78\% при +3 п.п.) --- редкие показы «слабой» руки дают мало информации. Отдельно проверено подглядывание за промежуточными результатами при фиксированном A/B: без поправки ошибка I рода растёт с 3\% до 11\%; процедуры OBF и mSPRT восстанавливают контроль (4\% и 2\%), но на логах бандита они IPS не заменяют. \emph{Итог:} для решения «B лучше A» нужен либо заранее запланированный A/B с фиксированным разбиением, либо IPS на уже собранных bandit-логах --- с оговоркой о мощности при слабом эффекте.

\textbf{Офлайн-оценка на исторических логах.} Реализован и проверен пайплайн IPS/SNIPS replay-OPE с корректной работой propensity на подвыборках OBD (E2, E2b/c, E9): на учебной выборке (10 тысяч событий) и на полном \texttt{random/all} (1,37 млн событий) метрики SNIPS, ESS и доля совпадений с логом считаются воспроизводимо. Однако надёжного ответа «какой алгоритм лучше» работа не получила: при низком CTR разница оценок SNIPS между политиками мала, а ранжирование сильно зависит от seed даже на миллионе событий. На логах адаптивного логирования (\texttt{bts/all}) эффективная выборка после взвешивания в разы ниже, чем на \texttt{random/all}. Проекция 80-ручного каталога на пару товаров (E13) наглядно показывает, как меняются метрики при сужении до двух рук, но 272 события не заменяют полноценный эксперимент. \emph{Итог:} OPE в курсовой --- рабочий инструмент для отладки пайплайна и осторожного сравнения кандидатов при равномерном логировании и достаточном числе кликов; при адаптивном логировании и скудных данных ранжирование политик по SNIPS продуктово не интерпретируется.

\textbf{Инженерная применимость.} Бенчмарк P1 при 80 руках показывает, что неконтекстные политики (Thompson Sampling, UCB, \texttt{fixed\_ab}) тратят на одно решение десятки микросекунд и не создают узкого места по времени отклика в типичном сервинговом контуре. LinUCB при том же числе рук требует порядка 23 ms на решение --- примерно в 450 раз дольше TS. Контекстный бандит оправдан только если ожидаемый прирост по кликам или персонализации (E8b) перекрывает этот вычислительный и инфраструктурный бюджет; в противном случае рациональнее неконтекстный TS или $\varepsilon$-greedy.

\textbf{Переносимость на другие метрики.} Все эксперименты проведены на бинарной награде «клик» --- в согласовании с OBD. При этом двухблочная схема (польза во время теста \emph{vs} корректная проверка гипотезы) и рекомендации табл. 4.1 не завязаны на CTR как на единственную бизнес-метрику: та же логика применима, если primary metric --- конверсия, выручка, удержание или другая наблюдаемая награда с достаточной частотой событий. Меняется калибровка (MDE, горизонт, мощность, редкость событий в OPE), а не принцип «MAB для накопления награды, отдельный A/B или IPS для вывода». Конкретные пороги в курсовой не проверялись --- это задача пилота или дипломной работы при появлении продуктовых логов.

\subsubsection{Научная новизна и вклад}\label{ux43dux430ux443ux447ux43dux430ux44f-ux43dux43eux432ux438ux437ux43dux430-ux438-ux432ux43aux43bux430ux434}

Теоретическое различие оптимизации и статистического вывода, риски наивного A/B на adaptive-логах и базовые OPE-процедуры известны \cite{ref5,ref13,ref19,ref21,ref26}. OBD \cite{ref26} и обзоры Burtini et al.~\cite{ref13}, Xiang et al.~\cite{ref21} не ставили задачу единого воспроизводимого фреймворка «польза онлайн-эксперимента → три типа сбоев статистического вывода → конкретная процедура исправления» в постановке CTR с открытым кодом.

Вклад курсовой --- методологический: (1) двухблочный воспроизводимый фреймворк в одном репозитории с тестами; (2) эмпирическая карта границ на синтетике и OBD --- где MAB выигрывает по regret, где даёт сбой статистический вывод, где OPE нестабилен; (3) количественная сверка IPS, OBF и mSPRT в единой постановке (E11, E14); (4) продуктовая матрица решений (табл. 4.1). Новые bandit-алгоритмы не предложены; работа систематизирует и проверяет известные тезисы для продуктовой линии CTR. Дополнительно собран пилотный воспроизводимый пайплайн: инфраструктура экспериментов (онлайн-оптимизация, статистический вывод, OPE, бенчмарк задержки) готова к переносу на внутренние продуктовые логи с propensity --- для обкатки рекомендаций из табл. 4.1 и принятия решения о раздельном или совмещённом дизайне MAB и A/B, когда такие данные появятся.

\subsection{Продуктовая рекомендация}\label{ux43fux440ux43eux434ux443ux43aux442ux43eux432ux430ux44f-ux440ux435ux43aux43eux43cux435ux43dux434ux430ux446ux438ux44f}

\textbf{Таблица 4.1 --- Ситуация → рекомендация → обоснование (эксперимент)}

Таблица сводит типичные продуктовые задачи к рекомендуемой процедуре. \emph{Ситуация} --- что нужно команде; \emph{Рекомендация} --- какой дизайн или способ анализа; \emph{Обоснование} --- соответствующий эксперимент курсовой.

\begingroup\footnotesize\setlength{\tabcolsep}{4pt}
\begin{xltabular}{\linewidth}{@{}
  >{\raggedright\arraybackslash}Y
  >{\raggedright\arraybackslash}Y
  >{\raggedright\arraybackslash}Y@{}}
\toprule\noalign{}
Ситуация & Рекомендация & Обоснование (эксп.) \\
\midrule\noalign{}
\endhead
\bottomrule\noalign{}
\endlastfoot
Нужно принять решение «внедрять B или нет» с контролем ошибки & Отдельный A/B с фиксированным разбиением & E4: ошибка I рода $\approx$ 3\%, мощность $\approx$ 85\% при T = 20k, +1 п.п. \\
Эксперимент уже шёл с бандитом, новый A/B нельзя & IPS-взвешенный статистический вывод по bandit-логам & E11: ошибка I рода 5\%; при слабом эффекте мощность низкая (24\%) \\
Смотрят промежуточные результаты фиксированного A/B & OBF или mSPRT (последовательный вывод) & E14: наивное подглядывание 11\% → OBF 4\%; не для bandit-логов \\
Важно не терять клики во время теста; много вариантов, быстрый отклик & MAB (TS, $\varepsilon$-greedy) & E1: regret меньше на 14--16; E12: +140 кликов \\
Отклик зависит от признаков пользователя & LinUCB, если допустимо время ответа & E8b: regret 26,6 vs 49,2; P1: $\sim$450× медленнее TS \\
Нужно сравнить политики без нового трафика & Офлайн-оценка по логам (replay-OPE, SNIPS); в логе должна быть propensity & E2/E9: пайплайн работает; устойчивого «рейтинга» политик нет \\
Логи собраны адаптивной политикой (BTS) & Не ранжировать политики по максимальному SNIPS & E2b/c, E9b: ESS в десятки раз ниже, чем на \texttt{random} \\
Мало событий и кликов (низкий CTR) & Не делать продуктовый вывод «кто лучше» по SNIPS & E2, E9: оценки близки, сильно зависят от seed \\
\end{xltabular}
\endgroup

\textbf{Таблица 4.2 --- Стоимость внедрения (относительная шкала, P1 + E1/E8b)}

\begingroup\footnotesize\setlength{\tabcolsep}{4pt}
\begin{xltabular}{\linewidth}{@{}
  >{\raggedright\arraybackslash}Y
  >{\raggedleft\arraybackslash}r
  >{\raggedleft\arraybackslash}r
  >{\raggedright\arraybackslash}Y
  >{\raggedright\arraybackslash}Y@{}}
\toprule\noalign{}
\midrule\noalign{}
\endhead
\bottomrule\noalign{}
\endlastfoot
\texttt{fixed\_ab} & 52,6 / 49,2 & 42 & низкая & для OPE --- да \\
\texttt{thompson\_sampling} & \textbf{36,6} / 28,3 & 51 & низкая & да \\
\texttt{ucb1} & 49,1 & 3 & низкая & да \\
\texttt{linucb} & --- / \textbf{26,6} & \textbf{23 080} & средняя & да (+ контекст) \\
\end{xltabular}
\endgroup

\textbf{Иллюстрация упущенной выгоды (E1, T = 5000).} Чтобы перевести regret в денежный порядок величины, взято учебное допущение: один клик условно стоит 100 ₽. Тогда разница regret между \texttt{fixed\_ab} и TS (около 16 кликов за эксперимент) соответствует примерно 1600 ₽ недобранной ценности на горизонте 5000 показов. Это не прогноз реальной выручки, а наглядная оценка: пока идёт эксперимент, адаптивное назначение может иметь заметную бизнес-стоимость --- параллельно с подготовкой статистически корректной проверки гипотезы.

\textbf{Ответ на центральный продуктовый вопрос.} Рекомендуемая практика --- не выбор между классическим A/B-тестом и бандитом, а связка подходов под разные цели. Курсовая количественно обосновывает обе части и предупреждает, где наивный вывод по логам бандита перестаёт быть корректным. Бандиты оправданы для оптимизации во время теста (блок I: измеримый выигрыш по кликам и regret). Статистический вывод по таким логам требует отдельных процедур --- фиксированный A/B (E4) или IPS-взвешенный анализ (E11); при слабом эффекте нужна доработка: больший горизонт, учёт propensity, развитие оценщиков и постановок анализа (DR, CUPED; направления в §4.3) --- как потенциал дипломной работы. Итог: работа показывает, как совместить оба шага --- накопление награды бандитом и проверку гипотезы --- без методологической ошибки.

\subsection{План экспериментов для развития работы до дипломной}\label{ux43fux43bux430ux43d-ux44dux43aux441ux43fux435ux440ux438ux43cux435ux43dux442ux43eux432-ux434ux43bux44f-ux440ux430ux437ux432ux438ux442ux438ux44f-ux440ux430ux431ux43eux442ux44b-ux434ux43e-ux434ux438ux43fux43bux43eux43cux43dux43eux439}

\subsubsection{Направление 1. Doubly Robust OPE vs SNIPS на OBD}\label{ux43dux430ux43fux440ux430ux432ux43bux435ux43dux438ux435-1.-doubly-robust-ope-vs-snips-ux43dux430-obd}

\begingroup\footnotesize\setlength{\tabcolsep}{4pt}
\begin{xltabular}{\linewidth}{@{}
  >{\raggedright\arraybackslash}Y
  >{\raggedright\arraybackslash}Y@{}}
\toprule\noalign{}
\midrule\noalign{}
\endhead
\bottomrule\noalign{}
\endlastfoot
\textbf{Гипотеза} & DR-оценка \cite{ref27} может снизить смещение и дисперсию SNIPS на adaptive-логах (BTS) при том же объёме данных \\
\textbf{Формулировка} & SNIPS\_DR стабильнее между seed на \texttt{bts/all}; ESS\_DR ≥ ESS\_SNIPS \\
\textbf{Метод} & Реализация DR; потоковый прогон на BTS-подвыборке OBD; 20 seed \\
\textbf{Критерии успеха} & Снижение std(SNIPS) между seed ≥ 20\%; документированный показатель стабильности ранжирования \\
\textbf{Риски} & Нестабильность DR при редких кликах; смягчение --- clipping весов, анализ чувствительности \\
\end{xltabular}
\endgroup

\subsubsection{Направление 2. CUPED + fixed A/B vs статистический вывод на bandit-логах}\label{ux43dux430ux43fux440ux430ux432ux43bux435ux43dux438ux435-2.-cuped-fixed-ab-vs-ux441ux442ux430ux442ux438ux441ux442ux438ux447ux435ux441ux43aux438ux439-ux432ux44bux432ux43eux434-ux43dux430-bandit-ux43bux43eux433ux430ux445}

\begingroup\footnotesize\setlength{\tabcolsep}{4pt}
\begin{xltabular}{\linewidth}{@{}
  >{\raggedright\arraybackslash}Y
  >{\raggedright\arraybackslash}Y@{}}
\toprule\noalign{}
\midrule\noalign{}
\endhead
\bottomrule\noalign{}
\endlastfoot
\textbf{Статус в курсовой} & Базовая реализация и E15 на синтетике выполнены; перенос на OBD --- для диплома \\
\textbf{Гипотеза} & CUPED \cite{ref7} при фиксированном разбиении может повысить мощность A/B; схема «MAB для exploration + отдельный CUPED-A/B для проверки гипотезы» потенциально предпочтительнее наивного статистического вывода на bandit-логах \\
\textbf{Формулировка} & При том же T мощность CUPED-fixed ≥ мощности IPS при слабом MDE (+0,5--1 п.п. CTR) \\
\textbf{Метод} & Синтетика + симуляция доэкспериментальной ковариаты; сравнение с E4/E11 \\
\textbf{Критерии успеха} & Мощность CUPED-fixed ≥ 90\% при +1 п.п. при T, где IPS ≤ 30\% \\
\textbf{Риски} & Нужны ковариаты; в OBD признаки ограничены --- сначала синтетика \\
\end{xltabular}
\endgroup

\subsubsection{Направление 3. Нестационарность CTR и деградация статистического вывода}\label{ux43dux430ux43fux440ux430ux432ux43bux435ux43dux438ux435-3.-ux43dux435ux441ux442ux430ux446ux438ux43eux43dux430ux440ux43dux43eux441ux442ux44c-ctr-ux438-ux434ux435ux433ux440ux430ux434ux430ux446ux438ux44f-ux441ux442ux430ux442ux438ux441ux442ux438ux447ux435ux441ux43aux43eux433ux43e-ux432ux44bux432ux43eux434ux430}

\begingroup\footnotesize\setlength{\tabcolsep}{4pt}
\begin{xltabular}{\linewidth}{@{}
  >{\raggedright\arraybackslash}Y
  >{\raggedright\arraybackslash}Y@{}}
\toprule\noalign{}
\midrule\noalign{}
\endhead
\bottomrule\noalign{}
\endlastfoot
\textbf{Гипотеза} & При дрейфе CTR \cite{ref16,ref38} regret-выигрыш MAB может сохраняться, а постфактумный статистический вывод --- деградировать быстрее, чем у fixed A/B с re-randomization \\
\textbf{Формулировка} & В нестационарной синтетике ошибка I рода IPS и наивного A/B растёт быстрее, чем в стационарных E4/E11 \\
\textbf{Метод} & Среда с дрейфом CTR; sliding-window IPS; сравнение с E1/E11 \\
\textbf{Критерии успеха} & Кривая ошибки I рода vs drift rate; рекомендации по re-randomization \\
\textbf{Риски} & Калибровка drift; начать с 2-arm \\
\end{xltabular}
\endgroup

\subsubsection{Направление 4. Продуктовый пилот (реальные логи или расширенный OBD)}\label{ux43dux430ux43fux440ux430ux432ux43bux435ux43dux438ux435-4.-ux43fux440ux43eux434ux443ux43aux442ux43eux432ux44bux439-ux43fux438ux43bux43eux442-ux440ux435ux430ux43bux44cux43dux44bux435-ux43bux43eux433ux438-ux438ux43bux438-ux440ux430ux441ux448ux438ux440ux435ux43dux43dux44bux439-obd}

\begingroup\footnotesize\setlength{\tabcolsep}{4pt}
\begin{xltabular}{\linewidth}{@{}
  >{\raggedright\arraybackslash}Y
  >{\raggedright\arraybackslash}Y@{}}
\toprule\noalign{}
\midrule\noalign{}
\endhead
\bottomrule\noalign{}
\endlastfoot
\textbf{Гипотеза} & Двухблочный фреймворк (§3.1) переносится на классическую A/B-постановку при достаточном трафике и согласии продукта; primary metric может быть не только CTR (конверсия, выручка и т.д.) при перекалибровке MDE и горизонта \\
\textbf{Формулировка} & На fixed A/B ошибка I рода контролируется; параллельный MAB даёт меньший regret при сопоставимом горизонте; для новой метрики --- отдельный прогон пайплайна с продуктовыми логами \\
\textbf{Метод} & Пилот или расширенная репликация на OBD; логирование propensity; отчёт по табл. 4.1 \\
\textbf{Экспертная оценка} & Полуструктурированные интервью с 3--5 продуктовыми менеджерами / аналитиками экспериментов: приемлем ли раздельный фреймворк «MAB для кликов + fixed A/B для проверки гипотезы»; шкала Лайкерта 1--5 по критериям из табл. 4.1 \\
\textbf{Критерии успеха} & Согласованный MDE; положительная экспертная оценка приемлемости (средний балл ≥ 4); воспроизводимый отчёт \\
\textbf{Риски} & Нет доступа к внутренним логам \cite{ref18}; смягчение --- углублённый OBD + синтетический якорь + экспертные интервью без production-данных \\
\end{xltabular}
\endgroup

\subsubsection{Направление 5. Strict contextual OPE для LinUCB}\label{ux43dux430ux43fux440ux430ux432ux43bux435ux43dux438ux435-5.-strict-contextual-ope-ux434ux43bux44f-linucb}

\begingroup\footnotesize\setlength{\tabcolsep}{4pt}
\begin{xltabular}{\linewidth}{@{}
  >{\raggedright\arraybackslash}Y
  >{\raggedright\arraybackslash}Y@{}}
\toprule\noalign{}
\midrule\noalign{}
\endhead
\bottomrule\noalign{}
\endlastfoot
\textbf{Гипотеза} & Strict contextual OPE \cite{ref6,ref28} подтверждает или опровергает выигрыш LinUCB из E8b на OBD (где batch-режим невалиден) \\
\textbf{Формулировка} & LinUCB ≥ TS по SNIPS на \texttt{random/all} с propensity и контекстом \\
\textbf{Метод} & Strict evaluator; subsample с контекстом; оптимизация времени задержки (P1) \\
\textbf{Критерии успеха} & SNIPS\_LinUCB \textgreater{} SNIPS\_TS с p \textless{} 0,05 bootstrap по seed; время задержки \textless{} 1 ms после оптимизации \\
\textbf{Риски} & Вычислительная стоимость; смягчение --- \texttt{np.linalg.solve}, sparse context \\
\end{xltabular}
\endgroup

\subsection{Ограничения, переносимость и воспроизводимость}\label{ux43eux433ux440ux430ux43dux438ux447ux435ux43dux438ux44f-ux43fux435ux440ux435ux43dux43eux441ux438ux43cux43eux441ux442ux44c-ux438-ux432ux43eux441ux43fux440ux43eux438ux437ux432ux43eux434ux438ux43cux43eux441ux442ux44c}

Работа ограничена одним доменом (fashion e-commerce, OBD) и одной эмпирически проверенной метрикой --- CTR (бинарный клик). Количественные выводы относятся к этой постановке; перенос на конверсию, выручку, удержание или непрерывные метрики отдельно не валидировался. На уровне процедур фреймворк метрико-агностичен (§1.8): награда задаётся в среде и в логе, а блоки I--II описывают роли дизайна, а не конкретный CTR --- см. абзац «Переносимость на другие метрики» выше. Monte Carlo-оценка статистического вывода выполнена на синтетике без промышленного A/B. DR \cite{ref27} описан в обзоре и §4.3, но не реализован. CUPED \cite{ref7} проверен в E15 на синтетике; на OBD не переносился. Нестационарность не проверялась (направление 3 в §4.3).

\textbf{Код и воспроизводимость.} Исходный код, тесты и скрипты экспериментов --- в репозитории курсовой; полный прогон --- \texttt{run\_all.ps1} / \texttt{run\_all.sh}. Сводные CSV для сверки с текстом --- в \texttt{outputs/} (см. \texttt{.gitignore}).

При наличии продуктовой платформы и данных каркас из §3.1 может быть адаптирован под пилот (§4.3, направление 4) без изменения логики блоков I и II.
